\documentclass[11pt]{article}
\usepackage[margin=1in]{geometry}
\usepackage{amsmath,amssymb}
\usepackage{booktabs}
\usepackage{array}
\newcolumntype{L}[1]{>{\raggedright\arraybackslash}p{#1}}
\usepackage{enumitem}
\usepackage[colorlinks=true,urlcolor=blue,linkcolor=black]{hyperref}
\setlist{itemsep=2pt,topsep=3pt}
\setlength{\parindent}{0pt}
\setlength{\parskip}{5pt}
\newcommand{\E}{\mathbb{E}}
\newcommand{\Var}{\operatorname{Var}}
\newcommand{\Cov}{\operatorname{cov}}

\title{Claims audit of the handwritten derivation pages:\\ the corrections, with derivations}
\author{Experiment 35, extended checks}
\date{11-26-2026}

\begin{document}
\maketitle

\section*{What this document shows and does not show}

\textbf{Shows.} Four of Brando's pages (the 12:03 and 13:01 derivation pages, and the annotated 6.437 page) contain 24 mathematical items. Each was checked with a Lean~4 proof over Mathlib (45 theorems, no \texttt{sorry}, only the standard axioms) and with a deterministic numeric test (36 of 36 pass). Of the 24, 12 are correct as written, 3 are correct with a note, 4 are wrong or cannot be shown, 2 are imprecise, 1 has a valid conclusion with an invalid reason, and 2 are questions. This PDF gives the corrected statements and their derivations. The full ledger, with the Lean theorem and numeric check behind each verdict, is \texttt{CLAIMS\_LEDGER\_11-26-2026.md} in the same folder.

\textbf{Does not show.} Nothing here concerns human raters or any trained judge: all instances are synthetic or exact. It does not test the claim that one loss or method is better than another; that is untested here. It does not re-prove the printed Theorem~1 of the 6.437 notes (cited). Three items are calculus facts checked numerically only (the pdf under rescaling, the Gaussian-noise reading of squared error, the integral form of the Brier divergence).

\textbf{What each objective can and cannot see.} This is the point behind items P3.3 and P3.4 below; $K$ is the number of ratings, $Z$ the input that fixes the pairing.

\begin{center}
\small
\begin{tabular}{@{}L{3.3cm}L{4.6cm}L{4.6cm}L{2.4cm}@{}}
\toprule
Objective & Sees & Does not see & Proper for the law? \\
\midrule
Mean difference $\E[X-Y]$ & $\mu_X-\mu_Y$ & shape, spread; $K-2$ free directions & no \\
Squared difference, fixed pairing & $\E[X\mid Z]$ & the rest of the conditional law & no \\
Squared difference, free pairing & whole law, after optimizing the pairing (squared Wasserstein) & for a given pairing, everything except $\mu_X,\mu_Y,\sigma_X,\sigma_Y,\rho$ & no \\
Squared difference, independent draws & $\sigma_p^2+\sigma_q^2+(\mu_p-\mu_q)^2$ & shape; rewards collapse to the mean & no (improper) \\
Correlation $\rho$ & linear association & any positive affine map of either variable & no \\
Brier score & whole pmf; excess $=\lVert p-q\rVert^2$ & nothing & yes (strictly); local only for $K=2$ \\
Log-loss & whole pmf & nothing & yes (strictly); local \\
\bottomrule
\end{tabular}
\end{center}

\section*{Notation}

$(X,Y)$ is a finite joint law; $\mu_X=\E[X]$, $\sigma_X^2=\Var[X]>0$, $\Cov(X,Y)=\E[(X-\mu_X)(Y-\mu_Y)]$, $\rho=\Cov(X,Y)/(\sigma_X\sigma_Y)$. As written on the pages, $\tilde X$ is the standardized variable. After the PS on page~1, $\tilde X=X-\mu_X$ is the centered variable and $Z_X=\tilde X/\sigma_X$ the standardized one; below I use the new convention. Ratings are $0,1,\dots,K-1$. A pmf $p$ and a report $q$ are vectors in the probability simplex.

\section*{Verdicts at a glance}

\begin{center}
\small
\begin{tabular}{@{}L{3.3cm}L{3.3cm}L{9.2cm}@{}}
\toprule
Item & Verdict & Line \\
\midrule
P1.1, P1.3--P1.6 & correct & top algebra; $\Cov$, $\mathrm{Corr}$, $Z$; $\Cov=\E[XY]-\E[X]\E[Y]$; $\Var=\E[X^2]-\mu^2$; the expansion of $\E[(X-Y)^2]$ and its boxed form \\
P1.2, S4 & questions & Q on Brier versus $\ell_2$; the margin questions on the 6.437 page \\
P1.7 & correct, with a note & boxed $2(1-\mathrm{corr})$ needs $Z$, not the centered tilde \\
P2.1, P2.3, P2.5, P2.6 & correct & $(a\pm b)^2$ lines; $\E[Z]=0$; $\Var$ forms; $\Var[Z]=1$ \\
\textbf{P2.2} & \textbf{wrong} & first expansion line misses the factor 2 \\
\textbf{P2.4} & \textbf{wrong} (struck out) & $\Var[\tilde X]=\Var[X-\mu]/\sqrt{\Var[\tilde X]}$ \\
P2.7 & imprecise & ``centering doesn't change the pmf/pdf'' \\
P2.8, P2.9 & correct, with a note & Brier divergence in pmf and pdf form; ``distribution level versus individual values'' \\
P3.1, P3.2 & correct & $X=Y\Leftrightarrow$ zero squared difference; the three-term expansion \\
\textbf{P3.3} & \textbf{wrong count} & ``6 degrees of freedom, 1 equation'': it is four \\
P3.4 & imprecise & ``can never do distribution matching'' holds for a fixed pairing only \\
\textbf{S1} & \textbf{cannot be shown} & ``show Brier is local'': false for $K\ge3$ \\
\textbf{S2} & reason invalid & ``log is 1-1 with its input'' does not give properness \\
S3 & correct & Bayes rule $\arg\min_a\E[C(a,H)\mid y]$ \\
\bottomrule
\end{tabular}
\end{center}

\section{P2.2 and P1.7: the factor two and the symbol change}

With $Z_X,Z_Y$ standardized, $\E[Z_X^2]=\E[Z_Y^2]=1$ and $\E[Z_XZ_Y]=\rho$. Expanding the square,
\[
\E[(Z_X-Z_Y)^2]=\E[Z_X^2]-2\,\E[Z_XZ_Y]+\E[Z_Y^2]=1-2\rho+1,
\qquad
\boxed{\E[(Z_X-Z_Y)^2]=2(1-\rho)}.
\]
The line written first on page 2 has the cross term $-Z_XZ_Y$ with no factor two. It evaluates to
\[
\E[Z_X^2-Z_XZ_Y+Z_Y^2]=1-\rho+1=2-\rho,
\]
which equals $2(1-\rho)$ only when $\rho=0$. Example: weights $(\tfrac14,\tfrac12,\tfrac14)$, $X=(0,1,2)$, $Y=(0,0,1)$ give $\mu_X=1$, $\sigma_X^2=\tfrac12$, $\mu_Y=\tfrac14$, $\sigma_Y^2=\tfrac3{16}$, $\Cov=\tfrac14$, so $\rho=\sqrt{2/3}\approx0.816$; the line as written gives about $1.184$, the correct value is about $0.367$.

If the tilde means the \emph{centered} variable (the PS), then $\E[(\tilde X-\tilde Y)^2]=\sigma_X^2+\sigma_Y^2-2\Cov(X,Y)$, which is $2(1-\rho)$ only if $\sigma_X=\sigma_Y=1$. So the boxed line must be written with $Z$, and the top-of-page definition becomes $Z_X=\tilde X/\sigma_X$. The raw-variable version, which needs no standardization, is
\[
\E[(X-Y)^2]=(\mu_X-\mu_Y)^2+\sigma_X^2+\sigma_Y^2-2\Cov(X,Y)
=(\mu_X-\mu_Y)^2+(\sigma_X-\sigma_Y)^2+2\sigma_X\sigma_Y(1-\rho).
\]

\section{P2.4: the variance of the standardized variable}

Dividing by $c=\sigma_X$ divides the variance by $c^2$, because $\Var[aX]=a^2\Var[X]$. Hence
\[
\Var[Z_X]=\E[Z_X^2]-\E[Z_X]^2=\frac{\E[(X-\mu_X)^2]}{\sigma_X^2}-0=\frac{\Var[X]}{\Var[X]}=1 .
\]
The scribbled line divides by $\sqrt{\Var[\tilde X]}$, which is the quantity being defined, and with the true value $\Var[\tilde X]=1$ its right side would equal $\Var[X]$, not $1$. It was right to cross it out.

\section{P2.7: what centering and standardizing do to a pmf and a pdf}

\emph{pmf.} If $T$ is injective, $P(T(X)=T(a))=P(X=a)$: every mass is carried to its new support point and none changes.

\emph{pdf.} For $Z=(X-\mu)/\sigma$, $F_Z(z)=F_X(\mu+\sigma z)$, so
\[
\boxed{p_Z(z)=\sigma\,p_X(\mu+\sigma z)} .
\]
Translation (centering) only moves the density; the division by $\sigma$ changes its height. For $X\sim\mathrm{Uniform}[0,1]$: $\mu=\tfrac12$, $\sigma=1/\sqrt{12}\approx0.2887$, and $p_Z=\sigma$ on $|z|<\sqrt3$ (area $2\sqrt3\,\sigma=1$), not the height $1$ of the original. A wording that holds: centering moves a pmf or pdf; standardizing also rescales a pdf's height; no such map changes a pmf's masses.

\section{P3.3: how many pmfs share a mean}

A signed measure $q\in\mathbb R^K$ on the ratings has total mass zero and its first moment zero when $\sum_i q_i\,i^k=0$ for $k=0,1$. In general take $k=0,\dots,j$: the $(j+1)\times K$ matrix $M_{ki}=i^k$ has rank $j+1$ whenever $j+1\le K$, since any $j+1$ of its columns form a Vandermonde matrix with distinct nodes and nonzero determinant. So the solution space has dimension
\[
\boxed{K-(j+1)} .
\]
For $K=6$: $j=1$ (normalization and mean) leaves $4$ free directions, $j=2$ leaves $3$, $j=5$ leaves $0$. The page's ``6 degrees of freedom, 1 equation'' should read: six unknowns, two equations (total mass and mean), four free directions. For $K=2$ the mean leaves $0$: it fixes the pmf. Four different pmfs on six ratings with mean $5/2$: uniform; $(\tfrac1{12},\tfrac16,\tfrac14,\tfrac14,\tfrac16,\tfrac1{12})$; $(\tfrac12,0,0,0,0,\tfrac12)$; $(0,0,\tfrac12,\tfrac12,0,0)$.

\section{P2.9 and P3.4: squared difference and distribution matching}

\textbf{Fixed pairing.} Let $Z$ fix the pairing and $f(Z)$ be the prediction of $X$. Then
\[
\E[(X-f(Z))^2]=\E[\Var(X\mid Z)]+\E\big[(\E[X\mid Z]-f(Z))^2\big],
\]
so the unique minimizer is $f=\E[X\mid Z]$ and the minimum is $\E[\Var(X\mid Z)]$. Two different conditional laws with the same conditional mean have the same minimizer: halves on $\{0,5\}$ and halves on $\{2,3\}$ both have mean $5/2$, with irreducible errors $25/4$ and $1/4$.

\textbf{Free pairing.} By Cauchy--Schwarz, $\Cov\le\sigma_X\sigma_Y$, so
\[
\boxed{\E[(X-Y)^2]\ \ge\ (\mu_X-\mu_Y)^2+(\sigma_X-\sigma_Y)^2},
\]
with equality when $\rho=1$. Zero squared difference forces $X=Y$ almost surely, hence equal laws. The minimum over all couplings of two given marginals is the squared Wasserstein distance, zero exactly for equal laws. Example: halves on $\{0,5\}$ against halves on $\{2,3\}$; coupling with $P(0,2)=P(5,3)=a$ costs $9-10a$ for $a\in[0,\tfrac12]$, so the minimum is $4$. So ``can never do distribution matching'' is true when the data fix the pairing, and false when the pairing is also optimized.

\textbf{Independent draws.} For $X\sim p$ and $Y\sim q$ independent,
\[
\boxed{\E[(X-Y)^2]=\sigma_p^2+\sigma_q^2+(\mu_p-\mu_q)^2},
\]
so the truthful report $q=p$ scores $2\sigma_p^2$ while a point mass at $y$ scores $\sigma_p^2+(y-\mu_p)^2$, which is smaller whenever some rating carrying mass is farther from $\mu_p$ than another rating is. For halves on $0$ and $5$: truthful $12.5$, point mass at $2$ or $3$: $6.5$. The rule is improper.

\section{S1: Brier is local only for two outcomes}

The Brier cost of a report $q$ at outcome $x$ is
\[
\sum_a\big(q_a-\delta_{a,x}\big)^2=\lVert q\rVert^2-2q_x+1,
\]
where $\delta_{a,x}=1$ if $a=x$ and $0$ otherwise.
For $K=2$, $q_0=1-q_1$ makes this $2(1-q_x)^2$, a function of $q_x$ alone: local. For $K=3$ it depends on the whole vector through $\lVert q\rVert^2$: $q_A=(\tfrac12,\tfrac12,0)$ and $q_B=(\tfrac12,\tfrac14,\tfrac14)$ both have $q_0=\tfrac12$, yet the cost at outcome $0$ is $\tfrac12$ against $\tfrac38$. So ``show Brier is local'' cannot be shown for $K\ge3$; this agrees with the printed Theorem~1 (log-loss is the only smooth, local, proper cost for $K\ge3$).

\section{S2: why log-loss is proper, and why ``one-to-one'' is not the reason}

A score of the form $\phi(q_x)$ has expected value $\sum_x p_x\,\phi(q_x)$ under the truth $p$, and is proper when this is maximized at $q=p$. Take $\phi(t)=t$, one-to-one in $t$, and $p=(\tfrac35,\tfrac25)$: reporting the truth earns $\tfrac9{25}+\tfrac4{25}=\tfrac{13}{25}$, while the point mass $(1,0)$ earns $\tfrac35=\tfrac{15}{25}$. So being one-to-one does not give properness ($\phi(t)=t^2$ fails the same way; for $\phi(t)=\sqrt t$ and $p=(0.6,0.4)$ the best report is $9/13\approx0.692$, not $0.6$).

For $\phi=\log$ the reason is Gibbs' inequality. Using $\log t\le t-1$,
\[
\sum_x p_x\log\frac{q_x}{p_x}\ \le\ \sum_x p_x\Big(\frac{q_x}{p_x}-1\Big)=\sum_x q_x-\sum_x p_x=0,
\qquad
\boxed{\sum_x p_x\log q_x\ \le\ \sum_x p_x\log p_x},
\]
with equality only if $q=p$.

\section*{Where the checks are}

All files are on \texttt{main} in the folder \nolinkurl{experiments/35_brier_vs_corr_math_l2_pearson_lean_start_26Nov2026/extended_checks_11-26-2026/claims_audit_11-26-2026/}:
\begin{itemize}
\item \texttt{AuditClaims.lean}: 45 theorems; \texttt{lean\_check.txt} holds the compile and axiom log.
\item \texttt{check\_claims\_audit.py}: 36 checks; output in \texttt{results/check\_claims\_audit.out}.
\item \texttt{CLAIMS\_LEDGER\_11-26-2026.md}: one entry per item, with the theorem name and check id.
\item \texttt{HOW\_TO\_READ\_RESULTS.md}: data, metrics and tolerances.
\end{itemize}

\end{document}
